\section{The appointed elements in context}
\subsection{Entrance: power invoked by the vulnerable}
Mordecai's prayer begins from a position of danger, not control. In Esther's court narrative, the people face destruction; appeal to the king must be made amid real peril. The prayer in the additions names God's power over everything and recalls creation. Its argument is that the empire's present threat does not exhaust reality. The visible authority capable of issuing a destructive command remains beneath the Creator's sovereignty.\footnote{Esther 4; 13:8--18; 14, Douay--Rheims Challoner. The Missal's reference cf. Esther 4:17 identifies the adapted Entrance in another numbering arrangement; Douay--Rheims 4:17 is not its text.}

This corporate danger matters to the antiphon's voice. Trust in divine power is not a claim to possess power over others. Mordecai intercedes for a people who cannot preserve themselves by willing security into existence. The creation language enlarges the horizon of petition: the same Lord who made heaven and earth can hear the people within them. The ensuing vineyard readings likewise concern a people whose existence begins in divine action. Yet the Entrance is an intercession amid danger, not an accusation that every endangered person deserves disaster.

The biblical prayer's continuation includes memory of deliverance and an appeal against destruction. Esther's own prayer adds the loneliness of one who must act within the court. Prayer and action therefore coexist in the larger narrative. The liturgical adaptation does not appoint every surrounding verse; those verses help explain why sovereignty is invoked. No securely checked direct patristic exposition of these precise Entrance verses supplies an additional identification here.

\subsection{Collect: petition under a greater generosity}
The Collect's generosity exceeds both the petitioner's merit and the reach of petition itself. Its three movements---divine abundance, mercy toward the guilty, and gifts beyond what can confidently be asked---make dependence more searching than a request for assistance with an otherwise self-sufficient life. The petitioner does not first produce an adequate claim and then receive its equivalent. God's bounty is the horizon within which the Church can ask.\footnote{Week 27 Collect, identified by \latin{Omnipotens sempiterne Deus, qui abundantia}; FDLC, \work{Mystagogical Reflections}, Ordinary Time 24--30 (2016), p. 12. The historical excerpt has an older conclusion; no approved prayer body is reproduced here.}

The petition for forgiveness also prevents generosity from becoming a refusal to name wrong. A person who needs mercy cannot make religious belonging a certificate of innocence. Beside Isaiah's disappointment and Matthew's violent tenants, the Collect gives the assembly a fitting posture: neither self-acquittal nor the despair that makes asking pointless. Augustine's account of being chosen so as to become good will deepen this relation in the first reading of the whole formulary. It is his account of John 15, rather than a commentary by him on this Collect.

\subsection{Isaiah: the love song becomes an indictment}
Isaiah 5 begins with the beloved and his vineyard. Its slope, cultivation, selected vine, tower and winepress establish attentive preparation. The poem invites its hearers to enter the situation sympathetically before it names them. The owner's question then asks Jerusalem and Judah to judge between him and the vineyard. The hearers who could readily recognize the injustice of disappointing such care discover that the case concerns their own life.\footnote{Isaiah 5:1--7; John Chrysostom, \work{Homilies on Matthew} 68.1--2, NPNF1, vol. 10.}

The final verse interprets the image: the vineyard is Israel and the chosen planting Judah. The expected fruit is judgment and justice; what arrives is iniquity and an outcry. The metaphor cannot be reduced to generic achievement, efficiency or successful institutions. Its fruit is the right ordering of life together. A vineyard can look substantial while the vulnerable cry out within it. The rest of Isaiah 5 makes the accusation socially concrete through its woes against acquisitiveness, self-indulgence and the corruption of judgment. Those woes are the chapter's explanatory context, not additional verses of the appointed reading.

The threatened removal of hedge and wall reverses the initial care. The place is to become exposed, uncultivated and dry. God has not been an indifferent owner discovering too late that his investment failed. The question is what more could have been done. The withdrawal of protection gives judgment a frightening form: a planting that refused its purpose loses the conditions within which it flourished. The passage holds divine love and accountability together, and its grief over failed fruit is part of its severity.

Chrysostom attends to precisely this image when explaining Matthew. He recognizes the prior care and the demand for works, but distinguishes the accused parties: Isaiah reproaches the vineyard, while Jesus' parable particularly reproaches the rulers responsible for it. That distinction changes the comparison. It allows the two texts to illuminate each other without identifying every vine with a violent tenant or assigning equal responsibility to everyone within the community.

\subsection{The psalm: the vine asks to live}
Psalm 80 in Hebrew numbering, 79 in the Vulgate and Douay--Rheims, recalls a vine brought out of Egypt and planted. The appointed verses move from that remembered gift to broken protection and devastation, then to petition for visitation and life. The assembly's answer is not merely assent to Isaiah's sentence. It becomes the prayer of those within the vineyard. Repeating the response from Isaiah 5:7a identifies the people, while the psalm lets that people lament.\footnote{Psalm 79:9,12--16,19--20, Douay--Rheims; Isaiah 5:7a. Augustine, \work{Exposition of Psalm} 79, §§6--11 (NPNF heading Psalm 80); Robert Bellarmine, \work{Commentary on the Book of Psalms}, Psalm 79, trans. John O'Sullivan (1866).}

Memory does work within the prayer. The vine once crossed from Egypt into a place prepared for it; its present wound is measured against that earlier care. Asking why the hedge is broken is therefore more than complaint about a bad season. It addresses the planter concerning his own work. The plea to look, see and visit seeks personal divine attention, while the appeal for life recognizes that even renewed invocation depends upon God's quickening.

The psalm also resists a simple historical illustration. The sea, river and ravaging beast invite identification, and the received commentators make different ones. Augustine understands the river as Jordan and the ravager through Roman power; Bellarmine understands the river as Euphrates and considers Assyrian and Babylonian devastation. Their Christian expositions agree more strongly on restoration than on the disaster's geography. Neither historical identification settles the composition date of the poem. The second interpretation develops their agreement and difference together.

\subsection{Philippians: prayer, attention and action}
Paul addresses a community, not an isolated mind learning to suppress discomfort. In the chapter's immediate context he calls Euodia and Syntyche to agreement in the Lord; afterward he speaks of the community's renewed material concern for him. Prayer, peace and practice belong within that web of common life. The appointed verses lead through three connected activities: bring requests with thanksgiving, attend to what is worthy, and do what has been received through apostolic teaching and example.\footnote{Philippians 4:2--3,6--9,10--20; Chrysostom, \work{Homilies on Philippians} 14, NPNF1, vol. 13. Verses 2--3 and 10--20 are context.}

Chrysostom begins his homily by asking how joy can coexist with Christ's blessing of those who mourn. His answer permits grief over sin and rejoicing in Christ together. Later, the peace surpassing understanding receives a theological explanation: God reconciles those who were hostile to him through the gift of his Son. Peace is not simply a sensation produced by choosing pleasant thoughts. It rests in a divine act whose generosity surpasses human expectation.

Paul's list then reaches conduct. Truth, purity, justice and what deserves praise are objects of attention because actions grow from thoughts. Chrysostom insists that the apostolic pattern includes what the hearers saw as well as what they heard: behavior teaches. He also distinguishes doing praiseworthy acts from acting for the sake of praise. The direction of the will matters as much as its public appearance. Applied beside Isaiah, this makes fruit something more demanding than a good reputation.

The exhortation against anxiety does not diagnose every anxious person as unfaithful. The text supplies a movement toward prayer, thanksgiving and practice within suffering; Chrysostom's own opening names affliction. His later strong language about the safety of virtue belongs to this moral exhortation. It cannot erase the suffering with which the homily begins or turn God's peace into a promise that no painful event will occur.

\subsection{The acclamation: chosen for lasting fruit}
John 15 belongs to Jesus' farewell discourse. The vine and branches, abiding, the command to love, and the giving of one's life for friends explain the fruit of which verse 16 speaks. The appointed acclamation is an adaptation identified by that verse, not an appointment of every clause of the biblical sentence. Its force within the Mass is to turn attention toward Christ's initiative and the enduring outcome of discipleship.\footnote{John 15:1--17; Augustine, \work{Tractates on John} 86.2--3 and 87.1, NPNF1, vol. 7.}

Augustine reads the order rigorously. Jesus does not choose people because they have already produced the fruit that makes them attractive candidates; he chooses them so that they may bear it. The next command, to love one another, names that fruit. The connection matters because ``fruit'' could otherwise become a blank into which ambition inserts its preferred success. Mutual charity gives the image a determinate Christian content. Augustine further orders petition in Christ's name toward salvation, excluding the idea that this name guarantees whatever desire happens to demand.

\subsection{The Gospel: entrusted land and a rejected Son}
Matthew 21 places the parable within Jerusalem's temple controversy. The leaders question Jesus' authority; the parables turn the question back upon their response to God. The vineyard is planted and equipped, then leased. The tenants' relation to it is therefore derivative from the beginning. When the owner sends for fruit, they answer with injury and murder. The sending of the son exposes their ambition most plainly: they try to seize inheritance by killing its heir.\footnote{Matthew 21:23--46, especially the appointed 33--43; Chrysostom, \work{Homilies on Matthew} 68.1--2.}

The parable's question brings a judgment from the hearers. They can identify the injustice when it is narrated about others. Jesus' citation of the rejected stone then discloses a divine reversal: the builders' assessment does not decide the stone's future. In the Christian reading the rejected Son is the cornerstone. The kingdom is given to a people bringing forth its fruits. The ending therefore answers seizure with responsibility, not with another group's entitlement to unaccountable possession.

Chrysostom finds providence in the prepared vineyard and longsuffering in repeated sending. He reads the owner's absence as patience and the expectation of respect for the son as the response that ought to have occurred, not ignorance in God. His homily includes fierce generalizations about Jews; these must be read alongside its own identification of the rulers and its statement that believing Jews and Gentiles become one. Matthew's following verses identify the chief priests and Pharisees as recognizing themselves in the parable. The accusation cannot honestly become hereditary blame attaching to every Jewish person.

The Church hears a parable about stewardship while receiving entrusted goods herself. Fruit remains required of the people to whom the vineyard is given. A reading that ends by congratulating its present hearers has lost the criterion by which the parable judges possession. The question is whether the owner receives the fruit of his gift, and whether the Son is received rather than displaced by the desire to rule.

\subsection{Offerings: obedience within sanctifying action}
The Prayer over the Offerings joins the Church's duty of worship to the sanctifying accomplishment of redemption. The inspected Latin witness supports that relation; an exact approved-English exemplar has not been established for this prayer. Its identifier and analysis therefore stand here without an English substitute prayer.\footnote{\latin{Suscipe, quaesumus, Domine}, Week 27; \work{Missale Romanum}, third typical edition (2002), secondary digital reproduction, physical pp. 290--291. This witness does not establish exact wording of the 2008 emended altar book or U.S. 2011 edition.}

The ritual place is significant. The Word's demand for fruit is followed by an offering that looks to God's sanctifying action. The worshipper does not leave divine gift behind in order to perform an independent act that forces a return. Aquinas explains Eucharistic efficacy from Christ and his Passion, then from the sacrament's nourishment and sign of unity. That theological account illuminates why an obedient offering and dependence on redemption belong together.\footnote{Thomas Aquinas, \work{Summa theologiae} III, q. 79, a. 1, body and replies.}

\subsection{Communion: waiting or incorporation}
Lamentations 3:25 speaks of the Lord's goodness to those who hope in him and seek him. The wider chapter has already spoken of lost peace and failed strength. Hope emerges within that speech of pain; it is not the report of someone untouched by ruin. The surrounding verses commend patience, but the chapter also rejects the crushing of prisoners and the perversion of justice. Waiting for God cannot be used to sanctify another person's oppression.\footnote{Lamentations 3:17--36,40--58; the appointed first Communion locus is 3:25.}

The alternative derived from First Corinthians 10:17 places common participation in the foreground. Verse 16's chalice and bread supply its scriptural setting; the Missal antiphon adapts the wording and includes chalice language. Paul's chapter warns against idolatry and closes by directing conduct toward others' good and God's glory. Communion is therefore a real belonging that makes rival allegiances and selfish conduct consequential. Chrysostom's explanation of many grains and one bread intensifies this claim: the same Christ feeds each person, so refusal of the neighbor contradicts the received communion.\footnote{1 Corinthians 10:14--24,31--33; Chrysostom, \work{Homilies on First Corinthians} 24.4, NPNF1, vol. 12.}

\subsection{After Communion: a life shaped by reception}
The Prayer after Communion seeks transformation through the nourishment received. The Church asks for an effect extending into life, rather than treating reception as the end of a transaction. Aquinas describes charity as aroused to action through the sacrament, and the body's members as serving justice now while awaiting future incorruption. His explanation joins moral action, sacramental gift and bodily hope.\footnote{Week 27, \latin{Concede nobis, omnipotens Deus}; FDLC, \work{Mystagogical Reflections}, p. 13; Aquinas, III, q. 79, a. 1, replies 2--3.}

Leo's Passion sermon supplies a further Eucharistic parallel: the receivers bear the crucified and risen Christ in body and spirit. The phrase about becoming what is received belongs to section VII's discussion of participation in Christ's body and blood; the preceding section discusses baptism. Leo's argument makes the Paschal identity of the gift explicit. Its resemblance to the final prayer supports theological understanding without supplying a newly established history of the modern prayer's composition.\footnote{Leo the Great, Sermon 63, \work{On the Passion XII}, VII and note 1063, trans. Charles Lett Feltoe, NPNF2, vol. 12.}
